\chapter{Ecosystem and Vocabulary}
\label{ch:ecosystem-vocabulary}
\label{ch:ecosystem-map}

CNA does not sit alone in the \texttt{openeggbert} GitHub organization. It sits at the
center of a small, deliberately layered set of sibling repositories, most of which expect to
be checked out as literal sibling directories on disk --- \texttt{../sharp-runtime},
\texttt{../easy-gl}, \texttt{../cna-samples}, and so on --- because CNake's own build system
looks for them there. This chapter is the map; later chapters (Part~VIII for
\texttt{sharp-runtime}, Part~IX for \texttt{easy-gl} and \texttt{free-direct}) are the
territory.

\section{The dependency graph}

\begin{center}
\begin{tabular}{lll}
\toprule
\textbf{Repository} & \textbf{Role} & \textbf{Consumed by} \\
\midrule
\texttt{sharp-runtime} & C\texttt{++} subset of the .NET BCL & CNA (foundation layer) \\
\texttt{easy-gl} & Toolkit-independent OpenGL/ES wrapper & CNA's five EasyGL-backed profiles \\
\texttt{free-direct} & DirectDraw/DirectSound/DirectPlay subset & CNA's \texttt{FREEDIRECT} renderer \\
\texttt{xna4-spec} & Machine-readable XNA 4.0 API spec (XML) & Compatibility audits of CNA \\
\texttt{cna-samples} & Ports of official XNA 4.0 samples & Exercises CNA end-to-end \\
\texttt{cna-examples} & In-app demo catalog for CNA & Exercises CNA end-to-end \\
\bottomrule
\end{tabular}
\end{center}

Every one of these is a genuinely independent, separately buildable, separately tested
project --- none of them is a vendored copy or a submodule of CNA. \texttt{sharp-runtime} in
particular describes itself as serving ``as a foundation for higher-level frameworks (e.g.
CNA)'' and lists CNA as its own downstream consumer, not the other way around; you could, in
principle, use \texttt{sharp-runtime} or \texttt{easy-gl} in a project that has nothing to do
with XNA at all.

\section{How the siblings are actually wired together at build time}
\label{sec:sibling-wiring}

The dependency graph above is not just a narrative claim --- it is exactly what CNA's own
\texttt{CMakeLists.txt} and \texttt{cmake/RendererSelection.cmake} do at configure time, checked
directly rather than assumed. Every sibling dependency follows the identical pattern: an
\texttt{if(NOT EXISTS \ldots/CMakeLists.txt)} guard that fails configuration outright, with a
real, copy-pasteable fix in the error message itself, rather than failing confusingly deep in a
missing-header compile error later:

\begin{lstlisting}[style=cnashell]
CNA: Missing sibling repository 'sharp-runtime' at
${CMAKE_CURRENT_SOURCE_DIR}/../sharp-runtime -- this is a separate git
checkout expected next to this repo's own directory, not a git
submodule (git submodule update --init will not fetch it). Fix:
cd ${CMAKE_CURRENT_SOURCE_DIR}/.. &&
git clone https://github.com/openeggbert/sharp-runtime.git
\end{lstlisting}

\texttt{sharp-runtime} is unconditionally required (\texttt{add\_subdirectory(../sharp-runtime
SHARP\_RUNTIME)} runs regardless of renderer choice); \texttt{easy-gl} and \texttt{free-direct}
are each pulled in only when a matching \texttt{-DCNA\_GRAPHICS\_RENDERER} value is selected,
so a plain \texttt{SDL\_RENDERER} or \texttt{VULKAN} build never even checks whether those two
sibling directories exist.

\begin{sourcenote}
\textbf{A real, second-level dependency this book's dependency-graph table doesn't show.}
\texttt{free-direct} is not a leaf in this graph --- its own \texttt{CMakeLists.txt} pulls in a
\emph{third} sibling repository, \texttt{free-api}, an experimental, small
WinAPI-compatibility-subset project (targeting roughly the Win32 API surface as it existed
around 1998) built on SDL3 the same way every other project in this ecosystem is, per
\texttt{free-api}'s own README. \texttt{RendererSelection.cmake}'s own comment explains why this
does not need a separate SDL configuration flag: \texttt{free-direct} resolves
\texttt{SDL3::SDL3} / \texttt{SDL3\_image::SDL3\_image} / \texttt{SDL3\_mixer::SDL3\_mixer} from
CNA's own already-vendored targets, set up before renderer selection runs, so the whole
three-level chain (CNA~$\to$~\texttt{free-direct}~$\to$~\texttt{free-api}) shares one SDL3
build rather than each level vendoring its own. The same comment names two real, independent
consumers of \texttt{free-direct} beyond CNA itself --- \texttt{free-eggbert} and
\texttt{planetblupi} --- confirming this sibling chain is genuinely reused ecosystem
infrastructure, not machinery built only for CNA's own \texttt{FREEDIRECT} renderer.
\end{sourcenote}

\section{sharp-runtime: the floor CNA stands on}
\label{sec:sharp-runtime-intro}

XNA's public API vocabulary assumes .NET underneath it: \texttt{byte}, \texttt{TimeSpan},
events, \texttt{IDisposable}, \texttt{System.Object.GetType()}. A C\texttt{++} reimplementation
of XNA has two choices --- invent ad-hoc C\texttt{++} substitutes for all of that inline,
scattered through the XNA porting work itself, or build the .NET-shaped foundation once, as
its own project, and depend on it. CNA's own \texttt{CLAUDE.md} porting rules choose the
second option explicitly: ``If CNA needs anything from .NET that is not yet in
\texttt{sharp-runtime} --- a type alias, a class, an interface, an exception type, a utility
--- add it to \texttt{sharp-runtime} first, then use it from CNA. Do not inline .NET concepts
directly into CNA headers as raw C\texttt{++} types or ad-hoc workarounds.''

\texttt{sharp-runtime} is explicit that it is \emph{not} attempting a CLR: no reflection, no
GC, no delegates beyond \texttt{std::function}, no serialization infrastructure, no P/Invoke,
no full cryptography/TLS stack. What it does provide is documented as a ``pragmatic subset
designed for use in native C\texttt{++} applications,'' tracked class-by-class in a
git-ignored local SQLite database (\texttt{plan.sqlite3}) with an honest, small state machine
per .NET type: \texttt{todo}, \texttt{ported}, \texttt{ignore} / \texttt{ignored}, or
\texttt{tobedecided} for genuinely ambiguous cases that need a human architecture call rather
than a guess. Chapter~\ref{ch:sharp-runtime-overview} and Chapter~\ref{ch:sharp-runtime-namespaces}
cover \texttt{sharp-runtime}'s own architecture, its \texttt{System::*} namespace layout, and its
component-isolated test apparatus in depth; Chapter~\ref{ch:parity-philosophy} covers the specific
question of how strictly CNA insists on that .NET shape versus where it deliberately diverges.

\section{easy-gl: OpenGL without a windowing opinion}

The EasyGL renderer family --- five public OpenGL/OpenGL ES profiles sharing one implementation
factory --- is built on \texttt{easy-gl}, a small,
toolkit-independent C\texttt{++}20 wrapper over OpenGL and OpenGLES. Its defining
architectural choice is what it refuses to own: no window creation, no GL context creation,
no event loop. The host application (CNA, in this case) owns the window, creates and
activates the GL context, and hands \texttt{easy-gl} a \texttt{GetProcAddress}-style loader
callback; \texttt{easy-gl} owns everything from there down --- shader compilation and
linking, buffers, vertex arrays, textures, draw calls, and capability queries that let calling
code gate OpenGL-only features (like polygon-mode fill flags) away from OpenGLES targets at
runtime via \texttt{device.supports(...)} / \texttt{device.require(...)}. As of the snapshot this
book draws from, \texttt{easy-gl}'s own README describes its scope as still a deliberately
compact, evolving vertical slice --- initialization, capability detection, and basic draw flow,
exercised by its own \texttt{hello-triangle-sdl} example and two dedicated smoke-test suites
(\texttt{easy-gl-smoke-tests}, \texttt{easy-gl-resource-smoke-tests}) --- with CNA's own
EasyGL implementation having grown well beyond that slice on top of it, per
Chapter~\ref{ch:easygl-backend}. Chapter~\ref{ch:easygl-deep-dive} covers \texttt{easy-gl}'s
public API surface (\texttt{easygl::Device}, \texttt{easygl::Shader}, \texttt{easygl::Program},
\texttt{easygl::Buffer}, \texttt{easygl::VertexArray}, \texttt{easygl::Texture},
\texttt{easygl::Feature}) in full.

\begin{sourcenote}
\textbf{A live instance of this book's own recurring pattern, caught in a sibling repo rather
than in CNA itself.} \texttt{easy-gl} carries its own \texttt{MIGRATION.md} (added 2026-07-19,
alongside a new 845-line \texttt{PLAN.md}), a nine-row ``breaking change summary'' table framed
entirely in future tense --- \texttt{Framebuffer::attach\_texture\_2d} ``now takes
\texttt{const~Texture\&}'', \texttt{Program::uniform\_block\_index} ``returns
\texttt{std::optional<unsigned~int>}'', and six more signature changes, each described as
``planned for the next release''. Reading the real, current headers directly instead of
trusting that framing shows every single one of the nine has already landed: checking a sample
of five directly against five real \texttt{include/easygl/} headers ---
\texttt{Framebuffer.hpp}, \texttt{Program.hpp}, and three others --- finds
\texttt{attach\_texture\_2d} is already taking
\texttt{const~Texture\&}, \texttt{uniform\_block\_index} already returning the
\texttt{optional}, and the remaining three checked signatures (\texttt{ProgramPipeline}'s two
\texttt{Program}-typed parameters, \texttt{TransformFeedback::set\_varyings},
\texttt{Sync::native\_handle()}'s \texttt{GLsync} return type) all already in their described
post-migration form. \texttt{MIGRATION.md} is not wrong about what changed --- it is simply
stale about \emph{when}, describing already-shipped API surface as still upcoming. This is the
identical shape of finding this book has made repeatedly about CNA's own tracked documents
(a snapshot that outlived the moment it was accurate), just observed here in a sibling
repository this book does not otherwise audit --- worth naming because it is independent
confirmation that the failure mode is a property of fast-moving source-controlled documentation
generally, not something specific to this project's own writing habits.
\end{sourcenote}

\section{free-direct: a narrow DirectX 3 for real legacy games}

Where \texttt{easy-gl} is a general-purpose OpenGL wrapper, \texttt{free-direct} is the
opposite kind of project: a deliberately \emph{narrow}, game-driven reimplementation of a
subset of DirectX~3 (2D) --- DirectDraw, DirectSound, and DirectPlay --- built on SDL3
internally, whose explicit goal is not DirectX compatibility in general but running two
specific legacy Windows games (\emph{Speedy Blupi} and \emph{Planet Blupi}) without their
original OS dependencies. Its own README states the non-goals as plainly as the goals: no
full DirectX~3 compatibility, no hardware-accurate emulation, no Direct3D pipeline, and no
expansion of APIs the target games do not actually call. This project backs CNA's
\texttt{FREEDIRECT} graphics renderer (Chapter~\ref{ch:dx3-freedirect-backend}) and is covered as a
project in its own right, including its real (not stubbed) ENet-backed DirectPlay networking,
in Chapter~\ref{ch:freedirect-deep-dive} and in the Speedy Blupi/Planet Blupi porting case study
in Chapter~\ref{ch:blupi-case-study}. \texttt{free-direct}'s own test suite is deliberately split
by transport: the default \texttt{ctest} run (no label filter) exercises
\texttt{LoopbackDirectPlayTransport}'s synchronous semantics and is not expected to pass
\texttt{directplay\_tests} unmodified against a real, asynchronous ENet build (1 of 8 tests fails
that way, by design, not as a regression); a separate \texttt{-L enet} label opts into the
ENet-backed transport tests specifically, verified 1/1 passing.

\section{xna4-spec: a ground truth to audit against}

A recurring failure mode for any reimplementation project is drift between what the README
claims and what the code actually does --- and a subtler one is drift between what the
README claims and what the \emph{original} API actually specified, from memory rather than
from a real reference. \texttt{xna4-spec} exists to remove the second failure mode. It is a
complete, machine-readable XML transcription of the official XNA 4.0 documentation --- every
class, struct, enum, interface, and delegate across all 19 XNA namespaces, 544 types in
total, with their original property/method/constructor/field/event descriptions --- converted
from Microsoft's own \texttt{learn.microsoft.com} XNA~4.0 reference pages. Its own README
states its purpose directly: ``Auditing C\texttt{++} reimplementations of XNA~4.0 (comparing
what is implemented vs.\ what is missing).'' The coverage figures quoted in this book's
Chapter~\ref{ch:what-is-cna} and throughout Parts~I--IV (227 of 245 public FNA types present,
92.7\%) are exactly this kind of audit --- a diff against a ground truth, not an estimate.
Chapter~\ref{ch:xna4-spec-auditing} walks through using \texttt{xna4-spec} to run this kind of
audit yourself, worked through in full against a real class (\cnaclass{Ray}) rather than
described only in the abstract.

\section{cna-samples and cna-examples: two different kinds of proof}

Two sibling repositories exist purely to give CNA something real to run, and they answer two
different questions.

\texttt{cna-samples} is a C\texttt{++} porting catalog for the official Microsoft XNA Game
Studio~4.0 sample archives, licensed Ms-PL to match its Microsoft-derived origin. Its question
is: \emph{does CNA run real, historical XNA sample code?}  At the pinned revision its plan
accounts for 153 archive entries: 63 completed ports, 23 tracked placeholders, and 67 deliberate
exclusions.  Those are catalog statuses, not fresh execution results; Chapter~\ref{ch:samples-and-examples}
explains why each dated placeholder must be re-audited before it is counted as a current CNA gap.

\texttt{cna-examples} is original, MIT-licensed CNA-specific work --- explicitly \emph{not}
a port of Microsoft's sample collection, precisely so it can carry a permissive license
\texttt{cna-samples} cannot. Conceptually inspired by JavaFX's \texttt{Ensemble} sample
browser, it is a single cross-platform application (desktop, Web, mobile) that lets a reader
navigate Area~$\to$~Category~$\to$~Demo (Input, Audio, Devices, Net, Media, 2D/3D Graphics,
\ldots) and run a live demonstration of each area, all inside one binary. Its question is:
\emph{does CNA have a coherent, in-app way to show off everything it can do?} As of the
snapshot this book draws from it is an early skeleton --- the application boots and
navigates Area~$\to$~Category, but individual demo screens are not yet implemented, and only
the EasyGL-backed renderer family is targeted so far.

\section{Repositories mentioned but out of this book's scope}

The \texttt{openeggbert} organization hosts a considerably larger set of repositories beyond
the six above --- additional CNA-adjacent projects (\texttt{cna-template}, \texttt{cna-extended},
\texttt{cna-craft}), other graphics libraries (\texttt{meta-gl}, \texttt{easy-3d}), actual
shipped games and tools built on this stack (\texttt{mesh-craft}, \texttt{galaxy-eggbert},
\texttt{free-eggbert}, \texttt{planetblupi}, \texttt{mobile-eggbert}), and a long tail of
per-project marketing/documentation sites (the \texttt{*.openeggbert.com} repositories). This
book's scope is deliberately the six repositories in the table above, because they are the
ones CNA's own build system, README, and \texttt{CLAUDE.md} porting rules directly name as
dependencies or verification partners. Appendix~\ref{ch:repo-map} gives a fuller repository
map for readers who want to go further.
